\section{第 12.1 节习题：验证任意方向的平移生成元}
\label{exercises:QM-C12-S01}
\exercisemeta
  {QM-C12-S01}
  {2026-10-17 至 2026-10-18}
  {Shankar, Second Edition；Exercise 12.1.1，印刷页 p.~306；本地 PDF 缺该页，题意据同版在线文字转录核对}
  {教材习题 12.1.1 已作答并核对}
  {笔记已确认；原页图像待补核}

\exercisequestion{QM-C12-S01-E01}{教材习题 12.1.1：从波函数平移识别生成元}
\textbf{对应知识点：}\relatedtopic{QM-C12-S01-T01}{方向上的动量分量}。题意统一记号重述，缺页及转录来源说明见\hyperref[section:QM-C12-S01]{正文}。

\textbf{题目。}设 $\mathbf n=(n_x,n_y)$ 是固定单位向量，$\psi(x,y)$ 足够光滑。主动平移位移 $\epsilon\mathbf n$ 后，波函数为
\begin{equation}
  \psi_\epsilon(x,y)=\psi(x-\epsilon n_x,y-\epsilon n_y).
\end{equation}
利用 $P_x=-\ii\hbar\partial_x$、$P_y=-\ii\hbar\partial_y$，将一阶展开与
$\psi_\epsilon=(I-\ii\epsilon G_{\mathbf n}/\hbar)\psi+O(\epsilon^2)$ 比较，验证沿 $\mathbf n$ 方向的无穷小平移生成元为 $\mathbf n\cdot\mathbf P$。

\begin{checkedsolution}
对两个自变量同时作一阶 Taylor 展开：
\begin{align}
  \psi_\epsilon(x,y)
  &=\psi(x,y)-\epsilon n_x\partial_x\psi(x,y)
    -\epsilon n_y\partial_y\psi(x,y)+O(\epsilon^2)\notag\\
  &=\left[I-\epsilon(n_x\partial_x+n_y\partial_y)\right]\psi
    +O(\epsilon^2)\notag\\
  &=\left[I-\frac{\ii\epsilon}{\hbar}(n_xP_x+n_yP_y)\right]\psi
    +O(\epsilon^2).
\end{align}
其中 $(-\ii)(-\ii)=-1$ 保证符号一致。对任意上述波函数比较生成元定义，得到
\begin{equation}
  \boxed{G_{\mathbf n}=n_xP_x+n_yP_y=\mathbf n\cdot\mathbf P}.
\end{equation}
\end{checkedsolution}

\section{第 12.2 节习题：旋转生成元与极坐标表示}
\label{exercises:QM-C12-S02}
\exercisemeta
  {QM-C12-S02}
  {2026-10-19}
  {Shankar, Second Edition；Exercise 12.2.2，印刷 p.~309 / PDF p.~317；Exercise 12.2.3，印刷 p.~310 / PDF p.~318}
  {12.2.2 采用教材允许的一致性验证要求；12.2.3 的两种方法均已作答并核对}
  {笔记已确认}

本题组只收录两道已完成的教材习题。12.2.1 涉及的逆自变量关系已在正文解释，未另设题；12.2.4 未选做，不收入课堂概念检查。

\exercisequestion{QM-C12-S02-E01}{教材习题 12.2.2：核对角动量的四条对易关系}
\textbf{对应知识点：}\relatedtopic{QM-C12-S02-T02}{分量变换与对易关系}。

\textbf{题目。}教材要求由旋转所需的对易关系确定 $L_z$，或至少验证 $L_z=XP_y-YP_x$ 与之相容。本题采用后一要求。已知
\begin{equation}
  [X,P_x]=[Y,P_y]=\ii\hbar I,\qquad
  [X,Y]=[P_x,P_y]=[X,P_y]=[Y,P_x]=0,
\end{equation}
验证
\begin{equation}
  [X,L_z]=-\ii\hbar Y,\quad [Y,L_z]=\ii\hbar X,\quad
  [P_x,L_z]=-\ii\hbar P_y,\quad [P_y,L_z]=\ii\hbar P_x.
\end{equation}

\begin{checkedsolution}
使用 $[A,BC]=[A,B]C+B[A,C]$。对位置分量，
\begin{align}
  [X,L_z]&=[X,X]P_y+X[X,P_y]-[X,Y]P_x-Y[X,P_x]\notag\\
         &=0+0-0-\ii\hbar Y=\boxed{-\ii\hbar Y},\\
  [Y,L_z]&=[Y,X]P_y+X[Y,P_y]-[Y,Y]P_x-Y[Y,P_x]\notag\\
         &=0+\ii\hbar X-0-0=\boxed{\ii\hbar X}.
\end{align}
对动量分量，注意对易子反对称性给出 $[P_x,X]=[P_y,Y]=-\ii\hbar I$，于是
\begin{align}
  [P_x,L_z]&=[P_x,X]P_y+X[P_x,P_y]-[P_x,Y]P_x-Y[P_x,P_x]\notag\\
           &=-\ii\hbar P_y+0-0-0=\boxed{-\ii\hbar P_y},\\
  [P_y,L_z]&=[P_y,X]P_y+X[P_y,P_y]-[P_y,Y]P_x-Y[P_y,P_x]\notag\\
           &=0+0-(-\ii\hbar P_x)-0=\boxed{\ii\hbar P_x}.
\end{align}
四条关系均成立。将它们用于 $U^\dagger A U=A+\ii\epsilon[L_z,A]/\hbar+O(\epsilon^2)$，位置和动量都按二维无穷小旋转规律变换。

\textbf{严格性补充。}这验证了一致性，却没有单凭对易关系唯一确定生成元：$L_z+cI$（$c$ 为实常数，具有角动量量纲）也满足同样的四条关系。它使旋转算符多出 $\ee^{-\ii\theta c/\hbar}$；正文选用不附加该整体相位的纯轨道旋转约定。
\end{checkedsolution}

\exercisequestion{QM-C12-S02-E02}{教材习题 12.2.3：用两种方法得到极坐标角动量}
\textbf{对应知识点：}\relatedtopic{QM-C12-S02-T02}{有限旋转与极坐标}。

\textbf{题目。}证明 $L_z=-\ii\hbar\partial_\varphi$，其中对 $\varphi$ 的偏导保持 $\rho$ 不变。要求分别使用：
\begin{enumerate}
  \item 坐标变换法：由 $L_z=-\ii\hbar(x\partial_y-y\partial_x)$ 和 $x=\rho\cos\varphi$、$y=\rho\sin\varphi$ 推出结论；
  \item 直接法：将主动小角度旋转 $\psi_\epsilon(\rho,\varphi)=\psi(\rho,\varphi-\epsilon)$ 与生成元定义 $U(\epsilon)=I-\ii\epsilon L_z/\hbar+O(\epsilon^2)$ 比较。
\end{enumerate}

\begin{checkedsolution}
\textbf{方法一：链式法则。}取任意足够光滑的 $f(x,y)$，定义
$F(\rho,\varphi)=f(\rho\cos\varphi,\rho\sin\varphi)$。固定 $\rho$ 时，
\begin{equation}
  \frac{\partial x}{\partial\varphi}=-\rho\sin\varphi=-y,\qquad
  \frac{\partial y}{\partial\varphi}=\rho\cos\varphi=x.
\end{equation}
于是
\begin{align}
  \frac{\partial F}{\partial\varphi}
  &=\frac{\partial f}{\partial x}\frac{\partial x}{\partial\varphi}
    +\frac{\partial f}{\partial y}\frac{\partial y}{\partial\varphi}\notag\\
  &=-y\frac{\partial f}{\partial x}+x\frac{\partial f}{\partial y}
   =(x\partial_y-y\partial_x)f.
\end{align}
因此，同一微分算符在极坐标表示中为
\begin{equation}
  \boxed{L_z=-\ii\hbar(x\partial_y-y\partial_x)
    =-\ii\hbar\partial_\varphi}.
\end{equation}

\textbf{方法二：无穷小旋转。}主动逆时针旋转后，当前极角处的振幅来自原来的 $\varphi-\epsilon$，所以
\begin{align}
  \psi_\epsilon(\rho,\varphi)
  &=\psi(\rho,\varphi-\epsilon)\notag\\
  &=\psi(\rho,\varphi)-\epsilon\partial_\varphi\psi(\rho,\varphi)
    +O(\epsilon^2).
\end{align}
另一方面，生成元定义给出
\begin{equation}
  \psi_\epsilon=\psi-\frac{\ii\epsilon}{\hbar}L_z\psi+O(\epsilon^2).
\end{equation}
对任意上述波函数比较一阶系数，
\begin{equation}
  -\partial_\varphi=-\frac{\ii}{\hbar}L_z
  \quad\Longrightarrow\quad
  \boxed{L_z=-\ii\hbar\partial_\varphi}.
\end{equation}
两种方法分别说明同一算符如何改变坐标表达，以及如何直接由几何变换识别生成元。角导数对应旋转，正如 $P_x=-\ii\hbar\partial_x$ 的位置导数对应平移。
\end{checkedsolution}
